% =====================================================================
\begin{thebibliography}{99}
\addcontentsline{toc}{chapter}{Bibliografía}
\markboth{Bibliografía}{}
\small

\bibitem{hammack}
R.~Hammack.
\newblock \emph{Book of Proof}, 3.ª ed. (edición~3.4).
\newblock Richard Hammack, 2018. Libro de texto de acceso libre.
\newblock \url{https://richardhammack.github.io/BookOfProof/}

\bibitem{mcs}
E.~Lehman, F.~T. Leighton y A.~R. Meyer.
\newblock \emph{Mathematics for Computer Science}, revisión del 18 de mayo de 2015.
\newblock Texto del curso 6.042J del MIT.
\newblock \url{https://ocw.mit.edu/courses/6-042j-mathematics-for-computer-science-spring-2015/}

\bibitem{mit6042}
A.~R. Meyer y A.~Chlipala.
\newblock 6.042J: Mathematics for Computer Science.
\newblock Massachusetts Institute of Technology, MIT OpenCourseWare, primavera de 2015. Unidad~1, «Proofs».
\newblock \url{https://ocw.mit.edu/courses/6-042j-mathematics-for-computer-science-spring-2015/}

\bibitem{polya}
G.~Polya.
\newblock \emph{How to Solve It: A New Aspect of Mathematical Method}.
\newblock Princeton University Press, 1945.

\bibitem{cs103}
Stanford University.
\newblock CS103: Mathematical Foundations of Computing.
\newblock Edición archivada \texttt{cs103.1272}, consultada en octubre de 2026: clases «Mathematical Proofs», «Indirect Proofs», «Propositional Logic» y «First-Order Logic»; guías «Guide to Proofs», «Guide to Negation» y «Proofwriting Checklist»; \emph{Problem Set}~1.
\newblock \url{https://web.stanford.edu/class/cs103/}

\bibitem{velleman}
D.~J. Velleman.
\newblock \emph{How to Prove It: A Structured Approach}, 3.ª ed.
\newblock Cambridge University Press, 2019.

\end{thebibliography}
